\newcommand{\CaseID}{river\_complex\_2021}
\newcommand{\CaseTitle}{2021 River Complex fire reconstruction against GOFER observations}
\newcommand{\FireName}{River Complex}
\newcommand{\FireYear}{2021}
\newcommand{\EventStart}{30 July 2021 at 21:00~UTC}
\newcommand{\EventEnd}{26 September 2021 at 13:00~UTC}
\newcommand{\DurationHours}{1384}
\newcommand{\CRSName}{UTM zone 10N (EPSG:32610)}
\newcommand{\FuelVersionText}{LANDFIRE release 2.2.0}
\newcommand{\FireSummary}{The supplied GOFER event record contains 1353 hourly perimeters, a final area of 931.308~km$^2$ and an official size of 199,359~ac. The configuration uses 5 timed ignition point(s).}
\newcommand{\ConfigurationText}{The supplied simulation runs for 1353~h, with an initial timestep of 1~s, target Courant--Friedrichs--Lewy number of 0.2, 1354 hourly weather fields and 5 ignition point(s). Arrival time, spread rate, fireline intensity and fire-size statistics are requested at the supplied dump interval. Crown-fire model 1 and firebrand spotting are enabled. Unspecified settings retain the inspected source defaults, including fixed random seed 2024, foliar moisture 90\% and maximum timestep 600~s. A single simulation member is configured.}
\newcommand{\WeatherOffsetText}{Simulation and weather time zero is 31~h after the GOES ignition hour; comparisons apply this offset explicitly.}
\newcommand{\ElmfireRevision}{8981f14e0170}
\newcommand{\ElmfireBranch}{yiren}
